\subsection{Affine Springer fibers} 

\subsubsection{Affine Springer fibers in the affine Grassmannian} For any $k$-algebra $R$, we denote $\frg\otimes_{k}R$ by $\frg(R)$. In particular, $\frg(F)=\frg\otimes_{k}F$ is the Lie algebra of the loop group $LG$. For $g\in LG$, let $\Ad(g)$ denote its adjoint action on $\frg(F)$. 

Let $\gamma\in\frg(F):=\frg\otimes_{k}F$ be a regular semisimple element \footnote{Here we are dealing with a Lie algebra $\frg$ over the non-algebraically-closed field $F$. An element  $\gamma\in\frg(F)$ is regular semisimple if it is regular semisimple as an element in $\frg(\overline{F})$, see the footnote in \S\ref{sss:Gr alt}. Equivalently, $\gamma$ is regular semisimple if its image in $\frc(F)$ lies in $\frc^{\rs}(F)$.}.  We consider the subfunctor of $\Gr_{G}$ whose value on a $k$-algebra $R$ is given by 
\begin{equation}\label{aff Spr}
\tX_{\gamma}(R)=\{[g]\in \Gr_{G}(R)|\Ad(g^{-1})\gamma\in\frg(R\tl{t})\}.
\end{equation} 
Then $\tX_{\gamma}$ is a closed sub-ind-scheme of $\Gr_{G}$. Let $\sX_{\gamma}=\tX^{\red}_{\gamma}$ be the underlying reduced ind-scheme of $\tX_{\gamma}$. We call $\sX_{\gamma}$ the {\em affine Springer fiber} of $\gamma$ in the affine Grassmannian $\Gr_{G}$.

\subsubsection{Alternative definition in terms of lattices}\label{sss:ASF Gr GLn} We consider the case $G=\GL_{n}$. Let $\gamma\in\frg(F)=\gl_{n}(F)$ be a regular semisimple matrix. As in \ref{sss:Gr} we identify $\Gr_{G}$ with the moduli space of lattices in $F^{n}$, or more precisely $\Gr_{G}(R)$ is the set of lattices in $R\lr{t}^{n}$. Then $\tX_{\gamma}(R)$ can be identified with those lattices $\L\subset R\lr{t}^{n}$ such that $\gamma \L\subset \L$, i.e., those stable under the endomorphism of $R\lr{t}^{n}$ given by $\gamma$. 

When $G=\SL_{n}$, $\Gr_{G}(R)$ classifies lattices $\L$ in $R\lr{t}^{n}$ such that $[\L:R\tl{t}^{n}]=0$. The affine Springer fiber $\tX_{\gamma}$ in this case is cut out by the same condition $\gamma \L\subset \L$.

When $G=\Sp_{2n}$, $\Gr_{G}(R)$ classifies lattices $\L$ in $R\lr{t}^{2n}$ such that $\L^{\vee}=\L$, see \S\ref{sss:Sp parahoric}. The affine Springer fiber $\tX_{\gamma}$ in this case is also cut out by the same condition $\gamma \L\subset \L$. 


We give the simplest examples of affine Springer fibers. 

\subsubsection{}\label{sss:reduction rs} Let $\gamma\in\frt(\calO_{F})$  such that the reduction $\ov\gamma\in\frt$ is regular semisimple. For each cocharacter $\l:\Gm\to T$, the element $t^{\l}:=\l(t)\in T(F)$ gives a point $[t^{\l}]\in \Gr_{G}(k)$ which lies in $\tX_{\gamma}$ since $\Ad(t^{-\l})\gamma=\gamma\in \frg(\calO_{F})$. The reduced ind-scheme $\sX_{\gamma}$ is in fact the discrete set $\{[t^{\l}]\}$ which is in bijection with $\xcoch(T)$. More canonically, there is an action of the loop group $LT$ on $\tX_{\gamma}$ given by its left translation action on $\Gr_{G}$. This action factors through the quotient $\Gr_{T}=LT/L^{+}T$ and realizes $\tX_{\gamma}$ as a $\Gr_{T}$-torsor.

\subsubsection{}\label{P1 chain} Consider the case $G=\SL_{2}$ and $\gamma=\mat{t}{0}{0}{-t}$. Then $\sX_{\gamma}$ is an infinite chain of $\PP^{1}$'s. More precisely, for each $n\in\ZZ$, we consider the subscheme $C_{n}$ of $\Gr_{G}$ classifying lattices $\L\subset F^{2}$ such that $t^{n}\calO_{F}\oplus t^{-n+1}\calO_{F} \subset \L\subset t^{n-1}\calO_{F}\oplus t^{-n}\calO_{F}$. Then $C_{n}\cong\PP^{1}$. We have $\sX_{\gamma}=\cup_{n\in\ZZ}C_{n}$ is an infinite chain of $\PP^{1}$'s. The components  $C_{n}$ and $C_{n+1}$ intersect at one point $t^{-n}\calO_{F}\oplus t^{n}\calO_{F}$ and otherwise the components are disjoint.

Here is a way to calculate the $k$-points of $\sX_{\gamma}$. We use the {\em Iwasawa decomposition} for $G(F)$: 
\begin{equation*}
G(F)=\bigsqcup_{n\in\ZZ}N(F)\mat{t^{n}}{0}{0}{t^{-n}}G(\calO_{F}).
\end{equation*}
According to this decomposition, any point in $\Gr_{G}$ can be represented by
\begin{equation}\label{Iwasawa elt}
g=\mat{1}{x}{0}{1}\mat{t^{n}}{0}{0}{t^{-n}}
\end{equation}
for some  $x\in F$ and a unique $n\in\ZZ$, and $x$ has a well-defined image in $F/t^{2n}\calO_{F}$. Since
\begin{equation*}
\Ad(g^{-1})\gamma=g^{-1}\gamma g=\mat{t}{2t^{1-2n}x}{0}{-t}
\end{equation*}
the condition $\Ad(g^{-1})\gamma\in\frg(\calO_{F})$ is the same as requiring $x\in t^{2n-1}\calO_{F}$. Therefore $\sX_{\gamma}(k)=\sqcup_{n\in\ZZ}Y_{n}$ where $Y_{n}$ consists of elements of the form \eqref{Iwasawa elt}  with $x\in t^{2n-1}\calO_{F}/t^{2n}\calO_{F}$. Therefore each $Y_{n}$ can be identified with $k$. It is easy to check that $Y_{n}\subset C_{n}(k)$.


\subsubsection{}\label{SL2 ram dim 1} Consider the case $G=\SL_{2}$ and $\gamma=\mat{0}{t^{2}}{t}{0}$. Then $\sX_{\gamma}$ consists exactly of those lattices $\L\in\Gr_{G}$ such that $t\calO_{F}\oplus \calO_{F}\subset\L\subset\calO_{F}\oplus t^{-1}\calO_{F}$. Therefore $\sX_{\gamma}\cong\PP^{1}$. Details of these calculations are left to the reader, see Exercises \ref{exII:SL2}.





\subsubsection{Invariance under conjugation} If $\gamma,\gamma'\in\frg(F)$ are related by $\Ad(g)\gamma=\gamma'$ for some $g\in G(F)$, then the left multiplication by $g$ on $\Gr_{G}$ restricts to an isomorphism $\tX_{\gamma}\cong\tX_{\gamma'}$, hence also $\sX_{\gamma}\cong\sX_{\gamma'}$. Therefore the isomorphism type of $\sX_{\gamma}$ is invariant under $G(F)$-conjugation on $\gamma$.  Recall we have map $\chi:\frg\to \frc:=\frg\sslash G\cong\frt\sslash W$. For a regular semisimple point $a\in\frc^{\rs}(F)$, the fiber $\chi^{-1}(a)$ is a single $G(F)$-conjugacy class (here we are using that the residue field of $F$ is algebraically closed). Therefore, the isomorphism type of $\sX_{\gamma}$  depends only on $a=\chi(\gamma)\in\frc^{\rs}(F)$. 

Unlike Springer fibers, $\sX_{\gamma}$ can be empty for certain $\gamma$. The affine Springer fiber $\sX_{\gamma}$ is nonempty if and only if $a=\chi(\gamma)\in\frc(\calO_{F})$. In fact, if $gG(\calO_{F})\in\sX_{\gamma}(k)$, then $\Ad(g^{-1})\gamma\in\frg(\calO_{F})$ hence $\chi(\gamma)=\chi(\Ad(g^{-1})\gamma)\in\frc(\calO_{F})$. Conversely, we have a Kostant section $\ep: \frc\to \frg$ of $\chi$ which identifies $\frc$ with $e+\frg_{f}$, where $(e,h,f)$ is a regular $\sl_{2}$-triple in $\frg$. Therefore, for any $a\in\frc(\calO_{F})\cap\frc^{\rs}(F)$, $\ep(a)\in\frg(\calO_{F})$, and $\sX_{\ep(a)}$ contains the unit coset in $\Gr_{G}$ hence nonempty; since $\sX_{\gamma}$ is isomorphic to $\sX_{\ep(\chi(\gamma))}$, it is also nonempty.

For $a\in\frc(\calO_{F})\cap\frc^{\rs}(F)$, we also write $\sX_{a}$ for $\sX_{\ep(a)}$. The above discussion shows that all Springer fibers $\sX_{\gamma}$ are isomorphic to $\sX_{a}$ for $a=\chi(\gamma)$.



\subsubsection{Parahoric versions} For each parahoric subgroup $\bP\subset LG$, we may similarly define the closed sub-indscheme $\tX_{\bP,\gamma}\subset\Fl_{\bP}$ using the analog of the condition \eqref{aff Spr} with $\frg(R\tl{t})$ replaced by $(\Lie\bP)\widehat{\otimes}_{k}R$. The reduced ind-scheme $\sX_{\bP,\gamma}=\tX_{\bP,\gamma}^{\red}\subset\Fl_{\bP}$ is called the {\em affine Springer fiber of $\gamma$ of type $\bP$}. In particular, when $\bI$ is an Iwahori subgroup of $LG$, we denote $\sX_{\bI,\gamma}$ by $\sY_{\gamma}$.

For $\bP\subset\bQ$ two parahoric subgroups of $LG$, the natural projection $\Fl_{\bP}\to \Fl_{\bQ}$ induces a map $\sX_{\bP,\gamma}\to \sX_{\bQ,\gamma}$. In particular we have a map $\sY_{\gamma}\to \sX_{\gamma}$.





\subsection{Symmetry on affine Springer fibers}
For the Springer fiber $\sB_{e}$, the centralizer $C_{G}(e)$ acts on it. In this subsection we investigate a similar structure for affine Springer fibers. 

\subsubsection{Centralizer action}\label{sss:centralizer} Let $G_{\gamma}$ be the centralizer of $\gamma$ in $G_{F}$ (the algebraic group over $F$ obtained from $G$ by base change). Then $G_{\gamma}$ is an algebraic group over $F$. Since $\gamma$ is regular semisimple, $G_{\gamma}$ is a torus over $F$. One can define the loop group $LG_{\gamma}$ of $G_{\gamma}$ as the functor $R\mapsto G_{\gamma}(R\lr{t})$ on $k$-algebras.

We claim that $LG_{\gamma}$ acts on the ind-scheme $\tX_{\gamma}$. This can be seen on the level of $R$-points. Suppose $h\in LG_{\gamma}(R)=G_{\gamma}(R\lr{t})$ and $[g]\in \tX_{\gamma}(R)$. Then the coset $[hg]\in \Gr_{G}(R)$ still satisfies
\begin{equation*}
\Ad((hg)^{-1})\gamma=\Ad(g^{-1})\Ad(h^{-1})\gamma=\Ad(g^{-1})\gamma\in\frg(R\tl{t})
\end{equation*}
using that $h$ is in the centralizer of $\gamma$. Therefore $[hg]\in \tX_{\gamma}(R)$. The assignment $[g]\mapsto [hg]$ for $h\in LG_{\gamma}$ and $[g]\in\tX_{\gamma}$ defines an action of   $LG_{\gamma}$ on $\tX_{\gamma}$. It induces an action of $LG_{\gamma}$ on the reduced structure $\sX_{\gamma}$.


\subsubsection{The split case} We consider the case where $\gamma\in\frt^{\rs}(F)$. In this case $G_{\gamma}=T\otimes_{k}F$, and
\begin{equation*}
LG_{\gamma}=LT\cong\xcoch(T)\otimes_{\ZZ}L\Gm
\end{equation*}
where $L\Gm$ is the loop group of the multiplicative group $\Gm$. For any $k$-algebra $R$, $L\Gm(R)=R\lr{t}^{\times}$. It is easy to see that an element $a=\sum_{i}a_{i}t^{i}\in R\lr{t}$ is invertible if and only if $a$ starts with finitely many nilpotent coefficients and the first non-nilpotent coefficient is invertible in $R$. When $R$ is reduced, the leading coefficient of $a$ must be invertible in $R$, which implies $R\lr{t}^{\times}=t^{\ZZ}\cdot R\tl{t}^{\times}$, and $R\tl{t}^{\times}=\{(a_{0},a_{1},\cdots)|a_{0}\in R^{\times}, a_{i}\in R, \forall i\geq 1\}$. We see that the reduced ind-scheme $(L\Gm)^{\red}\cong \ZZ\times L^{+}\Gm$, and that $L^{+}\Gm\cong \Gm\times\AA^{\NN}$ as schemes, where $\AA^{\NN}=\Spec k[x_{1},x_{2},\cdots]$. Therefore, when  $\gamma\in\frt^{\rs}(F)$, we have $(LT)^{\red}\cong\xcoch(T)\times L^{+}T$, and $L^{+}T$ is an affine scheme of infinite type.


\subsubsection{The lattice $\L_{\gamma}$}\label{sss:L gamma} For a general regular semisimple $\gamma\in\frg(F)$, let $\xcoch(G_{\gamma}):=\Hom_{F}(\Gm,G_{\gamma})$ be the $F$-rational cocharacter lattice of the torus $G_{\gamma}$. For each $\l\in\xcoch(G_{\gamma})$ viewed as a homomorphism $\Gm\to G_{\gamma}$ defined over $F$, we may consider the element $\l(t)$. The assignment $\l\mapsto \l(t)$ defines an injective homomorphism
\begin{equation*}
\xcoch(G_{\gamma})\incl G_{\gamma}(F).
\end{equation*}
whose image is denoted by $\L_{\gamma}$. It can be shown that the quotient $\L_{\gamma}\bs (LG_{\gamma})^{\red}$ is an affine scheme that is a finite disjoint union of $\Gm^{a}\times \AA^{\NN}$ for some integer $a$.



\subsubsection{The case $G=\GL_{n}$}\label{sss:GLn centralizer} We continue with the setup of \S\ref{sss:ASF Gr GLn}. We assume that $\chark>n$. Then the characteristic polynomial $P(x)=x^{n}+a_{1}x^{n-1}+\cdots+a_{n}\in F[x]$ of $\gamma$ is separable. The $F$-algebra $F[x]/(P(x))$ is then a product of fields $F_{1}\times \cdots \times F_{m}$, with $\sum_{i=1}^{m}[F_{i}:F]=n$. Each field extension $F_{i}/F$ is obtained by adjoining a root of an irreducible factor $P_{i}(x)$ of $P(x)$, and $F_{i}$ is necessarily of the form $k\lr{t^{1/e_{i}}}$ since $\chark>n$. Then the centralizer $G_{\gamma}$ is isomorphic to the product of {\em induced tori}
\begin{equation*}
G_{\gamma}\cong \prod_{i=1}^{m}\Res^{F_{i}}_{F}\Gm.
\end{equation*} 
We have $\xcoch(G_{\gamma})\cong \ZZ^{m}$, and the map $\xcoch(G_{\gamma})\to G_{\gamma}(F)$ is given by
\begin{equation*}
\ZZ^{m}\ni(d_{1},\cdots, d_{m})\mapsto (t^{d_{1}},\cdots, t^{d_{m}})\in F^{\times}_{1}\times\cdots\times F^{\times}_{m}
\end{equation*}
The quotient $\L_{\gamma}\bs G_{\gamma}(F)$ is isomorphic to $\prod_{i=1}^{m}F^{\times}_{i}/t^{\ZZ}$. Since each $F_{i}$ is isomorphic to $k\lr{t^{1/e_{i}}}$, we have  an exact sequence $1\to \calO^{\times}_{F_{i}}\to F^{\times}_{i}/t^{\ZZ}\to \ZZ/e_{i}\ZZ\to 0$, and hence the quotient $\L_{\gamma}\bs (LG_{\gamma})^{\red}$  contains the group scheme $\prod_{i=1}^{m}L^{+}_{F_{i}}\Gm$ with finite index. Here $L^{+}_{F_{i}}\Gm$ is isomorphic to $L^{+}\Gm$ as a scheme, except that we are renaming the uniformizer $t^{1/e_{i}}$. 

Alternatively, we may fix a uniformizer $t_{i}\in F_{i}$ (for example take $t_{i}=t^{1/e_{i}}$) and let $\wt\L_{\gamma}=t_{1}^{\ZZ}\times\cdots\times t^{\ZZ}_{m}\subset \prod_{i}F_{i}^{\times}=G_{\gamma}(F)$. The lattice $\wt\L_{\gamma}$ will be useful in calculating orbital integrals, see \S\ref{sss:GLn orb coho}.


\subsubsection{The case $G=\SL_{2}$} Let $G=\SL_{2}$ and $\gamma=\mat{0}{t^{n}}{1}{0}$ where $n\geq 1$ is odd. Then $G_{\gamma}(F)$ consists of matrices $\mat{a}{bt^{n}}{b}{a}$ with $a,b\in F$ and $a^{2}-t^{n}b^{2}=1$. Note that this equation forces $a,b\in\calO_{F}$, hence $G_{\gamma}(F)=G_{\gamma}(\calO_{F})$. The torus $G_{\gamma}$ is non-split and splits over the quadratic extension $E=F(t^{1/2})$. The lattice $\L_{\gamma}=\Hom_{F}(\Gm,G_{\gamma})=0$. Writing $a=\sum_{i\geq0}a_{i}t^{i}$ and $b_{i}=\sum_{i\geq 0}b_{i}t^{i}$, we see that $a_{0}=\pm1$, and once $b$ and $a_{0}$ are fixed, the higher coefficients of $a$ can be solved uniquely using the Taylor expansion of $(1+t^{n}b^{2})^{1/2}$. Therefore $(LG_{\gamma})^{\red}\cong L^{+}G_{\gamma}$ is isomorphic to $\{\pm1\}\times \AA^{\NN}$, given by $(a,b)\mapsto (a_{0},b_{0},b_{1},\cdots)$. 



\subsubsection{Symmetry on affine Springer fibers}\label{sss:sym ASF} Ng\^o has found a more precise statement about the action of $LG_{\gamma}$ on $\tX_{\gamma}$, namely the action factors through a canonical finite-dimensional quotient. We sketch the story following \cite[\S3.3]{NgoFL}. Let $a=\chi(\gamma)\in\frc(F)^{\rs}$ be the image of $\gamma$ under $\chi:\frg\to \frc$. We assume $a\in\frc(\calO_{F})$ for otherwise $\sX_{\gamma}$ is empty. 

There is a smooth affine group scheme $J$ over $\frc$ called the {\em regular centralizer group scheme}. It is characterized by the property that its pullback to $\frg$ via $\chi$, denoted $\chi^{*}J$, maps into the universal centralizer group scheme $I$ over $\frg$, and this map is an isomorphism over the regular locus $\frg^{\reg}$. Let $J_{a}$ be pullback of $J$ under the map $a:\Spec\calO_{F}\to \frc$. Then $J_{a}$ is a smooth affine group scheme over $\calO_{F}$ whose $F$-fiber is the torus $G_{\gamma}$ (i.e., $J_{a}$ is an integral model of $G_{\gamma}$ over $\calO_{F}$). We may form the positive loop group $L^{+}J_{a}$ of $J_{a}$ as well as its affine Grassmannian $P_{a}:=LG_{\gamma}/L^{+}J_{a}$ (also called the {\em local Picard group}). The reduced group scheme $P^{\red}_{a}$ is finite-dimensional and locally of finite type. Ng\^o showed that the action of $LG_{\gamma}$ on $\tX_{\gamma}$ (and hence on $\sX_{\gamma}$) factors through the local Picard group $P_{a}$, and it does not factor through any further quotient. For related statement, see \S\ref{dim formula discussion}.
 


